% ch15_representability_yoneda_philosophy.tex — Volume II Chapter 3

\chapter{Representability and the Yoneda Philosophy}

Volume I proved the Yoneda lemma and noted its importance. This chapter
takes the Yoneda perspective seriously as the central organizing principle
of all of category theory: \emph{to know an object is to know how everything
maps into and out of it}.

We prove the Yoneda lemma again, this time via the prooflog format that
makes the structure of the argument visible. Then we unpack the philosophical
content: universal properties are representability conditions, and the Yoneda
embedding is an embedding precisely because knowing hom-sets determines objects
up to isomorphism.

% ─────────────────────────────────────────────────────────────────────────────
\begin{nameorigin}
\term{Representable} (cross-ref Vol I Ch.~9 Yoneda) marks a functor
as ``representable by'' a specific object — analogous to representable
in classical algebra (a representation of a group acting on a space).
The \emph{Yoneda philosophy}: every important functor in mathematics is
representable, or close to it. The terminology was popularized by
Grothendieck's school in the 1960s; the modern emphasis on
representability as a structural principle is largely due to Mac Lane,
Lawvere, and Bénabou.
\end{nameorigin}

\section{Representable Functors}

\begin{definition}
A functor $F : \mathcal{C}^{op} \to \mathbf{Set}$ is \term{representable} if
there exists an object $c \in \mathcal{C}$ and a natural isomorphism
$F \cong \mathcal{C}(-, c)$.

Dually, a covariant functor $F : \mathcal{C} \to \mathbf{Set}$ is
\term{corepresentable} if $F \cong \mathcal{C}(c, -)$ for some $c$.

The object $c$ is called the \term{representing object}, and the natural
isomorphism is the \term{representation}.
\end{definition}

The Yoneda lemma says exactly how many representations an object has: natural
transformations $\mathcal{C}(-,c) \Rightarrow F$ are in bijection with elements
of $F(c)$. In particular, representations of $F$ by $c$ correspond to
``universal elements'' — special elements $u \in F(c)$ such that every
element of $F(d)$ is in the image of some morphism $d \to c$.

% ─────────────────────────────────────────────────────────────────────────────
\section{The Yoneda Lemma: Full Proof}

\begin{theorem}[Yoneda Lemma]
Let $\mathcal{C}$ be a locally small category, $c \in \mathcal{C}$, and
$F : \mathcal{C}^{op} \to \mathbf{Set}$. There is a bijection
\[
  \Phi : [\mathcal{C}^{op}, \mathbf{Set}](\mathcal{C}(-,c),\, F) \;\xrightarrow{\sim}\; F(c),
\]
natural in both $c$ and $F$, given by $\Phi(\alpha) = \alpha_c(\operatorname{id}_c)$.
\end{theorem}

\begin{prooflog}
\proofstep{Define $\Phi : \operatorname{Nat}(\mathcal{C}(-,c), F) \to F(c)$ by
           $\Phi(\alpha) = \alpha_c(\operatorname{id}_c)$.}
           {Evaluate the natural transformation $\alpha$ at $c$, then at the
            identity morphism $\operatorname{id}_c \in \mathcal{C}(c,c)$.}
\proofstep{Define $\Psi : F(c) \to \operatorname{Nat}(\mathcal{C}(-,c), F)$ as follows:
           for $x \in F(c)$, set $\Psi(x)_d(f) = F(f)(x)$ for $f : d \to c$.}
           {Given $x \in F(c)$ and a morphism $f : d \to c$, apply $F$ to $f$
            to get $F(f) : F(c) \to F(d)$, then evaluate at $x$.}
\proofstep{Check that $\Psi(x)$ is natural: for $g : e \to d$, we need
           $\Psi(x)_e(f \circ g) = F(g)(\Psi(x)_d(f))$, i.e.,
           $F(f \circ g)(x) = F(g)(F(f)(x))$.}
           {This holds because $F$ is a functor: $F(f \circ g) = F(g) \circ F(f)$.}
\proofstep{$\Phi \circ \Psi = \operatorname{id}$: $\Phi(\Psi(x)) = \Psi(x)_c(\operatorname{id}_c)
           = F(\operatorname{id}_c)(x) = x$.}
           {$F(\operatorname{id}_c) = \operatorname{id}_{F(c)}$ since $F$ is a functor.}
\proofstep{$\Psi \circ \Phi = \operatorname{id}$: given $\alpha$,
           $\Psi(\Phi(\alpha))_d(f) = F(f)(\alpha_c(\operatorname{id}_c))$.}
           {By definition of $\Psi$ and $\Phi$.}
\proofstep{Naturality of $\alpha$ at $f : d \to c$ says:
           $\alpha_d(f) = \alpha_d(\mathcal{C}(f,c)(\operatorname{id}_c))
           = F(f)(\alpha_c(\operatorname{id}_c))$.}
           {The naturality square for $\alpha$ at $f$: $\alpha_d \circ \mathcal{C}(f,c) = F(f) \circ \alpha_c$.}
\proofstep{So $\Psi(\Phi(\alpha))_d(f) = F(f)(\alpha_c(\operatorname{id}_c)) = \alpha_d(f)$.
           Hence $\Psi(\Phi(\alpha)) = \alpha$.}
           {This holds for all $d$ and all $f : d \to c$, confirming $\Psi \circ \Phi = \operatorname{id}$.}
\proofstep{The bijection $\Phi$ is natural in $c$: for $g : c \to c'$,
           $\Phi(\mathcal{C}(-,g) \circ \alpha) = \alpha_{c'}(g)$; and
           $F(g)(\Phi(\alpha)) = F(g)(\alpha_c(\operatorname{id}_c)) = \alpha_{c'}(g)$.}
           {Functoriality and the naturality of $\alpha$ at $g$.}
\end{prooflog}

% ─────────────────────────────────────────────────────────────────────────────
\section{The Yoneda Embedding}

\begin{theorem}[Yoneda Embedding]
The functor $Y : \mathcal{C} \to [\mathcal{C}^{op}, \mathbf{Set}]$ defined by
$Y(c) = \mathcal{C}(-, c)$ and $Y(f) = \mathcal{C}(-, f)$ is fully faithful.
\end{theorem}

\begin{prooflog}
\proofstep{We need to show: for all $c, c' \in \mathcal{C}$, the function
           $Y_{c,c'} : \mathcal{C}(c,c') \to [\mathcal{C}^{op},\mathbf{Set}](Y(c), Y(c'))$
           is a bijection.}
           {Fully faithful = bijective on all hom-sets.}
\proofstep{Apply Yoneda with $F = Y(c') = \mathcal{C}(-,c')$:
           $[\mathcal{C}^{op},\mathbf{Set}](\mathcal{C}(-,c),\,\mathcal{C}(-,c'))
           \cong \mathcal{C}(-,c')(c) = \mathcal{C}(c,c')$.}
           {The Yoneda bijection with $F = Y(c')$.}
\proofstep{Under this bijection, $f : c \to c'$ corresponds to the natural
           transformation $\mathcal{C}(-,f)$ given by post-composition with $f$.}
           {Tracing through $\Psi$: $\Psi(f)_d(g) = \mathcal{C}(d,f)(g) = f \circ g = \mathcal{C}(-,f)_d(g)$.}
\proofstep{So $Y_{c,c'}$ is the Yoneda bijection, hence a bijection. $Y$ is fully faithful.}
           {The Yoneda lemma directly implies the Yoneda embedding is fully faithful.}
\end{prooflog}

\begin{dialog}
``So the Yoneda embedding says that $\mathcal{C}$ embeds into its presheaf category
as the full subcategory of representable presheaves?''

Exactly. The Yoneda embedding is fully faithful, meaning $\mathcal{C}$ and the
subcategory of representable presheaves in $[\mathcal{C}^{op}, \mathbf{Set}]$ are
equivalent. Every category can be ``linearized'' into a presheaf category without
losing any information.
\end{dialog}

% ─────────────────────────────────────────────────────────────────────────────
\section{Universal Properties as Representability}

The Yoneda philosophy has a practical payoff: every universal property is a
representability statement, and the Yoneda lemma converts between them.

A product $a \times b$ in $\mathcal{C}$ is an object representing the functor
\[
  T_{a,b}(c) = \mathcal{C}(c, a) \times \mathcal{C}(c, b) : \mathcal{C}^{op} \to \mathbf{Set}.
\]
The universal pair of projections $(\pi_1, \pi_2)$ is the universal element in
$T_{a,b}(a \times b) = \mathcal{C}(a\times b, a) \times \mathcal{C}(a\times b, b)$.

A coequalizer of $f, g : a \to b$ is an object representing
\[
  E_{f,g}(c) = \{h : b \to c \mid h \circ f = h \circ g\} \subseteq \mathcal{C}(b, c).
\]

An exponential $b^a$ in a cartesian closed category represents
$c \mapsto \mathcal{C}(c \times a, b)$, and the evaluation morphism is the
universal element.

\begin{theorem}
An object $r \in \mathcal{C}$ together with a natural isomorphism
$F \cong \mathcal{C}(-,r)$ is precisely an initial object of the category of
elements $\int F$ of $F$.
\end{theorem}

The \term{category of elements} $\int F$ of $F : \mathcal{C}^{op} \to \mathbf{Set}$
has objects $(c, x)$ with $x \in F(c)$ and morphisms $(c,x) \to (c', x')$ the
morphisms $f : c' \to c$ in $\mathcal{C}$ (note the reversal) with $F(f)(x) = x'$.
The universal element is initial in this category.

% ─────────────────────────────────────────────────────────────────────────────
\section{Exercises}

\begin{exercisesbox}
\begin{qa}
\qaitem{What does it mean for $F : \mathcal{C}^{op} \to \mathbf{Set}$ to be representable?}
       {There exists $c \in \mathcal{C}$ with $F \cong \mathcal{C}(-,c)$ naturally.}

\qaitem{What is the representing object?}
       {The object $c$ such that $F \cong \mathcal{C}(-,c)$.}

\qaitem{State the Yoneda lemma as a bijection.}
       {$[\mathcal{C}^{op}, \mathbf{Set}](\mathcal{C}(-,c), F) \cong F(c)$,
        natural in $c$ and $F$.}

\qaitem{What is the bijection map $\Phi$ in the Yoneda lemma?}
       {$\Phi(\alpha) = \alpha_c(\operatorname{id}_c)$: evaluate $\alpha$ at $c$
        then at the identity.}

\qaitem{What is the inverse $\Psi$ of $\Phi$?}
       {$\Psi(x)_d(f) = F(f)(x)$ for $f : d \to c$ and $x \in F(c)$.}

\qaitem{What key step proves $\Psi \circ \Phi = \operatorname{id}$?}
       {The naturality of $\alpha$: $\alpha_d(f) = F(f)(\alpha_c(\operatorname{id}_c))$.}

\qaitem{What does it mean for the Yoneda embedding to be fully faithful?}
       {The function $\mathcal{C}(c,c') \to \operatorname{Nat}(\mathcal{C}(-,c),\mathcal{C}(-,c'))$
        is a bijection for all $c, c'$.}

\qaitem{Is the Yoneda embedding an equivalence of categories?}
       {Not in general: it need not be essentially surjective (not every presheaf
        is representable).}

\qaitem{What is the Yoneda philosophy in one sentence?}
       {An object is completely determined by its hom-functors — knowing
        $\mathcal{C}(-,c)$ or $\mathcal{C}(c,-)$ determines $c$ up to isomorphism.}

\qaitem{What functor does a product $a \times b$ represent?}
       {$T_{a,b}(c) = \mathcal{C}(c,a) \times \mathcal{C}(c,b)$: pairs of
        morphisms from $c$ to both $a$ and $b$.}

\qaitem{What is a universal element of a representable functor $F$?}
       {An element $u \in F(r)$ corresponding to $\operatorname{id}_r$ under
        the Yoneda bijection; equivalently, the image of $\operatorname{id}_r$
        under the representation isomorphism.}

\qaitem{What is the category of elements $\int F$ of $F : \mathcal{C}^{op} \to \mathbf{Set}$?}
       {Objects: pairs $(c, x)$ with $x \in F(c)$. Morphisms $(c,x) \to (c',x')$:
        morphisms $f : c' \to c$ in $\mathcal{C}$ with $F(f)(x) = x'$.}

\qaitem{What does the universal element correspond to in $\int F$?}
       {The initial object of $\int F$.}

\qaitem{What does the Yoneda lemma say about natural transformations between
        representable functors?}
       {$\operatorname{Nat}(\mathcal{C}(-,c), \mathcal{C}(-,c')) \cong \mathcal{C}(c,c')$:
        natural transformations between representables are exactly morphisms
        between the representing objects.}

\qaitem{Does the Yoneda lemma hold for enriched categories?}
       {Yes, with the appropriate enriched version: $[\mathcal{C},\mathcal{V}]_\mathcal{V}
        (\mathcal{C}(a,-), F) \cong F(a)$ in $\mathcal{V}$.}

\qaitem{What does it mean that ``the Yoneda embedding is the free cocompletion''?}
       {$[\mathcal{C}^{op},\mathbf{Set}]$ is the universal cocomplete category receiving
        $\mathcal{C}$ via a functor that preserves all colimits that exist in $\mathcal{C}$.}

\qaitem{How does the Yoneda lemma imply that $\mathcal{C}$ determines
        $[\mathcal{C}^{op},\mathbf{Set}]$?}
       {The Yoneda embedding $Y : \mathcal{C} \hookrightarrow [\mathcal{C}^{op},\mathbf{Set}]$
        is fully faithful, so $\mathcal{C}$ embeds as a full subcategory, and
        $[\mathcal{C}^{op},\mathbf{Set}]$ is the free cocompletion of this subcategory.}

\qaitem{Why does $\mathcal{C}(-, c) \cong \mathcal{C}(-, c')$ imply $c \cong c'$?}
       {The Yoneda embedding is fully faithful, so it reflects isomorphisms.}

\qaitem{What is a contravariant representable presheaf?}
       {A functor $\mathcal{C}^{op} \to \mathbf{Set}$ of the form $\mathcal{C}(-,c)$
        for some $c$.}

\qaitem{Give the Yoneda lemma as a statement about adjunctions.}
       {The Yoneda embedding $Y : \mathcal{C} \to [\mathcal{C}^{op}, \mathbf{Set}]$
        gives a bijection: natural transformations $Y(c) \Rightarrow F$ are in
        natural bijection with elements of $F(c)$.}
\end{qa}
\end{exercisesbox}

\begin{summarybox}
A functor $F : \mathcal{C}^{op} \to \mathbf{Set}$ is \textbf{representable} when
$F \cong \mathcal{C}(-,c)$. The \textbf{Yoneda lemma} says
$\operatorname{Nat}(\mathcal{C}(-,c), F) \cong F(c)$ naturally, with bijection
$\alpha \mapsto \alpha_c(\operatorname{id}_c)$. The \textbf{Yoneda embedding}
$Y : \mathcal{C} \hookrightarrow [\mathcal{C}^{op}, \mathbf{Set}]$ is fully faithful,
embedding $\mathcal{C}$ into its presheaf category. Every universal property is a
representability condition; the universal element is the initial object of the
category of elements.
\end{summarybox}

\section{Formal Restatement}

\begin{formalrestatement}
\textbf{Representable.} $F : \mathcal{C}^{op} \to \mathbf{Set}$ representable
iff $\exists\, c,\, u \in F(c)$ with $\Psi(u) : \mathcal{C}(-,c) \xrightarrow{\sim} F$
a natural isomorphism.

\textbf{Yoneda lemma.}
$\operatorname{Nat}(\mathcal{C}(-,c), F) \cong F(c)$, naturally in $c \in \mathcal{C}$
and $F \in [\mathcal{C}^{op},\mathbf{Set}]$.
Bijection: $\Phi(\alpha) = \alpha_c(\operatorname{id}_c)$;
inverse: $\Psi(x)_d(f) = F(f)(x)$.

\textbf{Yoneda embedding.}
$Y : \mathcal{C} \hookrightarrow [\mathcal{C}^{op},\mathbf{Set}]$, $Y(c) = \mathcal{C}(-,c)$,
fully faithful. Consequence: $c \cong c'$ iff $Y(c) \cong Y(c')$;
$\mathcal{C}(-,c) \cong \mathcal{C}(-,c') \Rightarrow c \cong c'$.

\textbf{Category of elements.}
$\int F$: objects $(c,x)$ with $x \in F(c)$; morphisms $(c,x) \to (c',x')$
are $f : c' \to c$ with $F(f)(x) = x'$. Universal element = initial object of $\int F$.
\end{formalrestatement}
